\chapter{A Concrete Naming Certificate for $\DiagonalLimitRepresentativeUp$}
\label{ch:concrete-instances-diagonallimitrepresentative-namecert}
\origin{human}

The $\DiagonalLimitRepresentativeUp$ packet realizes the diagonal representative of \autoref{sec:real-completeness-regular-cauchy} as a finite naming certificate. It sits after the diagonal-modulus surface of \autoref{ch:concrete-instances-regular-cauchy-limit-modulus-namecert}, and it consumes the finite windows, readbacks, dyadic ledgers, and terminal real-name seal supplied by \autoref{ch:concrete-instances-streamname-namecert}, \autoref{ch:concrete-instances-regseqrat-namecert}, \autoref{ch:concrete-instances-dyadicratcore-namecert}, and \autoref{ch:concrete-instances-real-namecert}. The same terminal real-facing handoff is the Cauchy-limit seal boundary of \autoref{ch:concrete-instances-cauchylimitseal-namecert}; the diagonal representative consumes it only as a displayed finite route, not as quotient-stream equality, a selected host limit, choice, or ambient Real completeness. The packet names the representative row itself; it is not a Real-completeness principle and it does not choose a host limit.

\begin{definition}[Diagonal-limit representative carrier]
\label{def:diagonal-limit-representative-carrier}
A $\DiagonalLimitRepresentativeUp$ carrier is a finite $\BHist$ packet
$$
D=(X,\mu,n,k,i,w,r,d,e,H,C,P,N)
$$
with the following displayed rows:
\begin{enumerate}
\item $X$ is the regular-Cauchy family over $\RegSeqRatUp$ whose diagonal representative is being inspected;
\item $\mu$ is the explicit modulus row read through \autoref{ch:concrete-instances-regular-cauchy-limit-modulus-namecert};
\item $n$ is the displayed precision request;
\item $k=\mu(1/(n+1))$ is the selected family index;
\item $i=k+n+1$ is the selected readback index;
\item $w$ is the $\StreamNameUp$ window exposing the $i$-th read of the $k$-th family member;
\item $r$ is the $\RegSeqRatUp$ readback row for that selected member and window;
\item $d$ is the $\DyadicRatCoreUp$ ledger recording the tolerance and rational arithmetic of the readback;
\item $e$ is the $\RealUp$ seal row reached after the selected regular-rational readback has been repacked;
\item $H$ records componentwise $\hsame$ transport for $X,\mu,n,k,i,w,r,d,e$ and the local certificate row;
\item $C$ records displayed $\Cont$ replay from the modulus request through the window, readback, dyadic ledger, and real seal;
\item $P$ is the $\Pkg$ provenance over $\DiagonalLimitRepresentativeUp$, $\RegSeqRatUp$, $\StreamNameUp$, $\DyadicRatCoreUp$, $\RealUp$, $\BHist$, $\hsame$, $\Cont$, and $\NameCert$;
\item $N$ is the local $\NameCert_{\DiagonalLimitRepresentativeUp}$ row.
\end{enumerate}
The carrier exports no host real, quotient stream equality, hidden tail selector, countable choice row, or ambient $\RealUp$ completeness field.
\end{definition}

\begin{definition}[Diagonal-limit representative classifier]
\label{def:diagonal-limit-representative-classifier}
Two accepted $\DiagonalLimitRepresentativeUp$ carriers are classified together exactly when their modulus rows select the same displayed diagonal route and their selected family index, readback index, $\StreamNameUp$ window, $\RegSeqRatUp$ readback, $\DyadicRatCoreUp$ ledger, $\RealUp$ seal row, transport rows, continuation rows, package rows, and local naming rows are related componentwise by the recorded $\hsame$ comparisons and $\Cont$ routes.
\end{definition}

\begin{theorem}[Diagonal-limit representative NameCert obligations]
\label{thm:diagonal-limit-representative-namecert-obligations}
For the carrier of \autoref{def:diagonal-limit-representative-carrier}, the local $\NameCert_{\DiagonalLimitRepresentativeUp}$ surface consists exactly of carrier admission, diagonal-index computation $k=\mu(1/(n+1))$ and $i=k+n+1$, finite-window exposure through $w$, regular-rational readback through $r$, dyadic ledger exposure through $d$, terminal seal exposure through $e$, componentwise transport through $H$, continuation replay through $C$, package provenance through $P$, and local naming through $N$.
\end{theorem}
\begin{proof}
Project the packet $D$. The carrier lists the family row, modulus row, request row, computed index rows, selected window, regular-rational readback, dyadic ledger, real seal, transport, continuation, provenance, and local naming row.

The diagonal-index computation is finite. The modulus surface supplies $k$, the unary request supplies $n$, and the readback index $i$ is obtained by the displayed rational-index arithmetic used in \autoref{def:real-completeness-diagonal-limit}.

The remaining rows expose the route through the four concrete sibling surfaces: $\StreamNameUp$ for $w$, $\RegSeqRatUp$ for $r$, $\DyadicRatCoreUp$ for $d$, and $\RealUp$ for $e$. Since $H,C,P,N$ only transport, replay, source, and name those displayed rows, the listed obligations exhaust the local certificate.
\end{proof}
\leanchecked{BEDC.Derived.DiagonallimitrepresentativeUp.DiagonalLimitRepresentativeNameCertObligations}

\begin{theorem}[Diagonal-limit representative route predicate obligation]
\label{thm:diagonal-limit-representative-route-predicate-obligation}
For an accepted $\DiagonalLimitRepresentativeUp$ carrier
$$
D=(X,\mu,n,k,i,w,r,d,e,H,C,P,N),
$$
the accepted route predicates are exactly the displayed rows $X,\mu,n,k,i,w,r,d,e,H,C,P,N$. The family and readback rows are read through $\RegSeqRatUp$, the selected window through $\StreamNameUp$, the dyadic ledger through $\DyadicRatCoreUp$, and the terminal seal through $\RealUp$. The support predicates are precisely componentwise $\hsame$ transport through $H$, $\Cont$ replay through $C$, $\Pkg$ provenance through $P$, and the local $\NameCert_{\DiagonalLimitRepresentativeUp}$ row $N$.
\end{theorem}
\begin{proof}
Project the carrier of \autoref{def:diagonal-limit-representative-carrier}. Its rows already name the family, modulus, request, selected indices, selected window, readback, dyadic ledger, Real seal, and the four support coordinates.

Apply the NameCert obligations of \autoref{thm:diagonal-limit-representative-namecert-obligations}. They certify the two index computations and expose the sibling rows in the order $\StreamNameUp$, $\RegSeqRatUp$, $\DyadicRatCoreUp$, and $\RealUp$.

It remains to exclude additional route predicates. The classifier of \autoref{def:diagonal-limit-representative-classifier} allows transport only by componentwise $\hsame$ and replay only by the recorded $\Cont$ routes, while $P$ and $N$ source and name the same packet. Hence the accepted route predicate surface is exactly the displayed carrier row.
\end{proof}

\begin{theorem}[Diagonal-limit representative terminal stability obligation]
\label{thm:diagonal-limit-representative-terminal-stability-obligation}
Let
$$
D=(X,\mu,n,k,i,w,r,d,e,H,C,P,N)
$$
and
$$
D'=(X',\mu',n',k',i',w',r',d',e',H',C',P',N')
$$
be accepted $\DiagonalLimitRepresentativeUp$ carriers classified together by their recorded componentwise $\hsame$ rows and replayed $\Cont$ routes. Then the selected diagonal indices, selected $\StreamNameUp$ window, $\RegSeqRatUp$ readback, $\DyadicRatCoreUp$ ledger, and terminal $\RealUp$ seal remain on the displayed route
$$
\mu,n,k,i,w,r,d,e.
$$
The transport does not introduce a hidden selector, detached window, alternate readback, second dyadic ledger, or off-packet terminal Real seal.
\end{theorem}
\begin{proof}
First read index stability. By \autoref{thm:diagonal-limit-representative-index-stability}, the displayed equations for $k,i$ and $k',i'$ are part of the accepted route and remain tied to the modulus request under classifier transport.

Next read the terminal route. The window-ledger exhaustion theorem, \autoref{thm:diagonal-limit-representative-window-ledger-exhaustion}, confines public reads of $w,r,d,e$ to the displayed sibling rows, and \autoref{thm:diagonal-limit-representative-real-seal-factorization} places every Real-facing read downstream of those rows.

Finally use the support rows. Componentwise $\hsame$ and replayed $\Cont$ preserve the rows already present in the packet; $\Pkg$ and $\NameCert$ name that same packet. Therefore terminal stability is stability of the displayed route, not permission to manufacture a hidden selector or an off-packet terminal seal.
\end{proof}

\begin{theorem}[Diagonal-limit representative obligation-closure package]
\label{thm:diagonal-limit-representative-obligation-closure-package}
The chapter's $\DiagonalLimitRepresentativeUp$ obligation surface is the package consisting of the carrier of \autoref{def:diagonal-limit-representative-carrier}, the classifier of \autoref{def:diagonal-limit-representative-classifier}, the NameCert obligations of \autoref{thm:diagonal-limit-representative-namecert-obligations}, the route predicate obligation of \autoref{thm:diagonal-limit-representative-route-predicate-obligation}, and the terminal stability obligation of \autoref{thm:diagonal-limit-representative-terminal-stability-obligation}. This package is the concrete diagonal representative surface available to later Lean formalization; it contains no quotient stream equality, host limit, selected completed limit, countable choice, or ambient $\RealUp$ completeness principle.
\end{theorem}
\begin{proof}
The carrier and classifier supply the accepted packet and its componentwise comparison rule. The NameCert obligations certify the displayed diagonal-index computations, sibling rows, and support rows.

The route predicate obligation identifies the exact predicate surface consumed by downstream readers. The terminal stability obligation then proves that componentwise transport and continuation replay preserve the selected diagonal route through the terminal Real seal.

Combining these pieces yields the chapter's obligation surface. The exclusions are inherited from the carrier and from the already established nonchoice and completion boundaries, so the package closes the obligation level without adding quotient, host-limit, choice, or ambient-completeness data.
\end{proof}

\begin{theorem}[Diagonal-limit representative primitive-scope exhaustion]
\label{thm:diagonal-limit-representative-primitive-scope-exhaustion}
For every accepted $\DiagonalLimitRepresentativeUp$ carrier
$$
D=(X,\mu,n,k,i,w,r,d,e,H,C,P,N),
$$
the displayed route predicates are scoped exactly by $\BHist$, $\hsame$, $\Cont$, $\Pkg$, $\NameCert$, and the sibling rows supplied by $\StreamNameUp$, $\RegSeqRatUp$, $\DyadicRatCoreUp$, $\RealUp$, and $\FinitePrefixLimitBudgetUp$. The representative route uses $X,\mu,n,k,i,w,r,d,e,H,C,P,N$ together with the shared-terminal finite-prefix rows of \autoref{def:diagonal-limit-representative-shared-terminal-carrier}; it has no primitive outside those listed BEDC rows.
\end{theorem}
\begin{proof}
Project the carrier. The only local rows are the family, modulus, request, selected indices, window, readback, ledger, seal, transport, continuation, provenance, and naming rows listed in \autoref{def:diagonal-limit-representative-carrier}.

Read the sibling surfaces. The route predicate obligation places $w,r,d,e$ respectively in $\StreamNameUp$, $\RegSeqRatUp$, $\DyadicRatCoreUp$, and $\RealUp$, and the shared-terminal carrier adds only the finite-prefix budget rows from $\FinitePrefixLimitBudgetUp$.

The support rows are exactly $H,C,P,N$. Thus every transport, replay, source, and local name is accounted for by $\hsame$, $\Cont$, $\Pkg$, and $\NameCert$ over $\BHist$. Since the carrier and the shared-terminal extension export no further coordinates, the primitive scope is exhausted by the listed rows.
\end{proof}

\begin{theorem}[Diagonal-limit representative support-scope nonescape]
\label{thm:diagonal-limit-representative-hsame-cont-scope}
All accepted transport and replay for the $\DiagonalLimitRepresentativeUp$ route occurs through the recorded support rows $H,C,P,N$. No hidden selector, quotient stream, host limit, ambient completeness field, countable-choice row, detached window, alternate dyadic ledger, or off-packet terminal seal can enter the route through $\hsame$, $\Cont$, $\Pkg$, or $\NameCert$.
\end{theorem}
\begin{proof}
Begin with terminal stability, \autoref{thm:diagonal-limit-representative-terminal-stability-obligation}. It says that accepted componentwise transport preserves the displayed route $\mu,n,k,i,w,r,d,e$ and does not manufacture hidden selectors or off-packet terminal seals.

Use the nonchoice and completion boundaries, \autoref{thm:diagonal-limit-representative-nonchoice-boundary} and \autoref{thm:diagonal-limit-representative-completion-boundary}. These boundaries remove quotient stream equality, host limits, countable choice, and ambient completeness from the packet surface.

Finally read the support rows. The classifier permits only the componentwise $\hsame$ comparisons and $\Cont$ routes recorded by $H$ and $C$, while $P$ and $N$ source and name the same finite packet. Therefore accepted transport and replay remain inside the recorded support scope.
\end{proof}

\begin{theorem}[Diagonal-limit representative scoped-closure package]
\label{thm:diagonal-limit-representative-scoped-closure-package}
The $\DiagonalLimitRepresentativeUp$ scoped-closure surface is the package consisting of the obligation-closure package of \autoref{thm:diagonal-limit-representative-obligation-closure-package}, the primitive-scope exhaustion theorem of \autoref{thm:diagonal-limit-representative-primitive-scope-exhaustion}, and the support-scope nonescape theorem of \autoref{thm:diagonal-limit-representative-hsame-cont-scope}. This package binds the diagonal representative route to the displayed BEDC primitives without claiming Real completeness, quotient-stream equality, host limits, selected completed limits, countable choice, or machine verification.
\end{theorem}
\begin{proof}
The obligation-closure package supplies the accepted carrier, classifier, NameCert obligations, route predicate obligation, and terminal stability obligation. These theorems identify the diagonal representative route and the terminal Real seal boundary.

The primitive-scope exhaustion theorem identifies the primitive rows available to the route: $\BHist$, $\hsame$, $\Cont$, $\Pkg$, $\NameCert$, the four representative sibling surfaces, and the finite-prefix budget surface. The support-scope nonescape theorem confines every accepted transport and replay step to $H,C,P,N$.

Combining these two scope results with the obligation package gives a closed scoped surface for the paper chapter. The exclusions are precisely the absent coordinates named by the carrier, nonchoice boundary, completion boundary, and finite-prefix handoff, so the package records scoped closure without adding analytic completion data.
\end{proof}

\begin{closurestatus}{\DiagonalLimitRepresentativeUp}
  \constructivestory{The diagonal-limit representative is generated from $BHist$ rows that record diagonal indices, finite prefixes, window ledgers, stream-name to regular-rational handoffs, and the shared terminal route. The construction follows the familiar diagonal-limit idea, but every representative is read through the carrier rather than selected from an ambient completion. The $hsame$ relation transports index and terminal rows across equivalent histories, while $\Cont$ carries the finite-prefix and Real-seal handoffs forward. The chapter's boundary, uniqueness, selector-exclusion, completion-frontier, shared-terminal, and handoff theorems keep the representative bottom-up: a terminal diagonal representative exists only as the visible route assembled from the named finite ledgers.}
  \theoryclosure{\scopedClosure}
  \scopeclosed{This chapter binds the diagonal-limit representative carrier, classifier, NameCert obligations, the route predicate obligation of \autoref{thm:diagonal-limit-representative-route-predicate-obligation}, the terminal stability obligation of \autoref{thm:diagonal-limit-representative-terminal-stability-obligation}, the closure package of \autoref{thm:diagonal-limit-representative-obligation-closure-package}, primitive-scope exhaustion of \autoref{thm:diagonal-limit-representative-primitive-scope-exhaustion}, support-scope nonescape of \autoref{thm:diagonal-limit-representative-hsame-cont-scope}, scoped-closure packaging of \autoref{thm:diagonal-limit-representative-scoped-closure-package}, Real-seal and completion boundaries, uniqueness, selector-exclusion, shared-terminal route, and finite-prefix budget handoff.}
  \formalstatus{\axiomCleanV}
  \leantarget{BEDC.Derived.DiagonallimitrepresentativeUp.DiagonalLimitRepresentativeNameCertObligations}
  \bridgestatus{none}
  \origin{human}
  \notclaimed{It does not claim a selected completion point, quotient-stream equality, global limit uniqueness, or choice beyond the displayed diagonal route.}
  \upgradepath{A higher closure grade requires Lean verification and a public bridge for the scoped package.}
\end{closurestatus}

\begin{theorem}[Diagonal-limit representative Real-seal factorization]
\label{thm:diagonal-limit-representative-real-seal-factorization}
Every $\RealUp$ consumer read of an accepted $\DiagonalLimitRepresentativeUp$ carrier factors through the ordered finite route
$$
\mu,n,k,i,w,r,d,e,H,C,P,N.
$$
The consumer reaches the seal row $e$ only after the modulus-selected family index, the selected $\StreamNameUp$ window, the $\RegSeqRatUp$ readback, and the $\DyadicRatCoreUp$ ledger have been exposed on the same packet.
\end{theorem}
\begin{proof}
Begin with the NameCert obligations. They force every accepted carrier to compute $k$ and $i$ before any readback row is visible.

Next read the sibling surfaces in order. The $\StreamNameUp$ window $w$ exposes the finite observation window, the $\RegSeqRatUp$ row $r$ supplies the selected regular-rational readback, and the $\DyadicRatCoreUp$ ledger $d$ records the tolerance and arithmetic data attached to that readback.

The $\RealUp$ seal row $e$ is downstream of these rows and is preserved only through the displayed $H,C,P,N$ support. Therefore a consumer cannot reach $e$ through a host real, quotient stream equality, detached tail selector, countable choice row, or ambient completeness field.
\end{proof}

\begin{theorem}[Diagonal-limit representative nonchoice boundary]
\label{thm:diagonal-limit-representative-nonchoice-boundary}
An accepted $\DiagonalLimitRepresentativeUp$ packet exports exactly the displayed family row $X$, modulus row $\mu$, request row $n$, selected indices $k,i$, selected $\StreamNameUp$ window $w$, $\RegSeqRatUp$ readback $r$, $\DyadicRatCoreUp$ ledger $d$, $\RealUp$ seal row $e$, componentwise transport $H$, continuation replay $C$, package provenance $P$, and local naming row $N$. It exports no host real, quotient stream equality, hidden tail selector, countable choice row, or ambient $\RealUp$ completeness.
\end{theorem}
\begin{proof}
Inspect the carrier coordinates of \autoref{def:diagonal-limit-representative-carrier}. Every coordinate is either a displayed BEDC row, a computed finite index, or structural support for transport, replay, provenance, and naming.

\autoref{thm:diagonal-limit-representative-namecert-obligations} accounts for the admissible carrier and classifier rows. \autoref{thm:diagonal-limit-representative-real-seal-factorization} accounts for every Real-facing consumer read of the packet.

The excluded objects require additional coordinates absent from the carrier. Hence they cannot be exported by the local naming certificate or recovered by a consumer of the displayed diagonal-limit representative route.
\end{proof}

\begin{theorem}[Diagonal-limit representative index stability]
\label{thm:diagonal-limit-representative-index-stability}
Let $D=(X,\mu,n,k,i,w,r,d,e,H,C,P,N)$ and $D'=(X',\mu',n',k',i',w',r',d',e',H',C',P',N')$ be accepted $\DiagonalLimitRepresentativeUp$ carriers classified together by the componentwise $\hsame$ rows recorded in their transport data. Then the displayed index equations are preserved:
$$
k=\mu(1/(n+1)),\qquad i=k+n+1,
$$
and
$$
k'=\mu'(1/(n'+1)),\qquad i'=k'+n'+1.
$$
Consequently the selected family index and selected readback index remain the same diagonal route under accepted componentwise transport, and no classifier read may replace them by a hidden selector.
\end{theorem}
\begin{proof}
Project both carriers through \autoref{def:diagonal-limit-representative-carrier}. The rows $k,i$ and $k',i'$ are not auxiliary choices; they are the displayed computations tied to $\mu,n$ and $\mu',n'$ inside the packet.

Use the classifier of \autoref{def:diagonal-limit-representative-classifier}. Classification requires the modulus rows, selected family indices, readback indices, windows, readbacks, ledgers, seals, transports, continuations, provenance, and local naming rows to be related componentwise by the recorded $\hsame$ comparisons and $\Cont$ routes.

Finally apply the NameCert obligations of \autoref{thm:diagonal-limit-representative-namecert-obligations}. Since the local certificate includes the two index computations themselves, transport of the packet preserves those displayed computations rather than producing a new tail selector. The accepted diagonal route is therefore stable under the componentwise classifier.
\end{proof}

\begin{theorem}[Diagonal-limit representative window-ledger exhaustion]
\label{thm:diagonal-limit-representative-window-ledger-exhaustion}
Every public read of the window, readback, ledger, and seal coordinates of an accepted $\DiagonalLimitRepresentativeUp$ carrier
$$
D=(X,\mu,n,k,i,w,r,d,e,H,C,P,N)
$$
factors through the displayed sibling rows $w,r,d,e$. The row $w$ is read through $\StreamNameUp$, $r$ through $\RegSeqRatUp$, $d$ through $\DyadicRatCoreUp$, and $e$ through $\RealUp$, with $H,C,P,N$ only transporting, replaying, sourcing, and naming those same rows.
\end{theorem}
\begin{proof}
Start with the carrier row. The only window, regular-rational readback, dyadic arithmetic ledger, and Real-facing seal coordinates are $w,r,d,e$.

Read the public route in the order forced by \autoref{thm:diagonal-limit-representative-real-seal-factorization}. The selected window $w$ exposes the finite $\StreamNameUp$ observation, the readback $r$ exposes the $\RegSeqRatUp$ row selected from that window, the ledger $d$ records the $\DyadicRatCoreUp$ tolerance and arithmetic data, and the seal $e$ is the downstream $\RealUp$ row.

It remains to exclude extra public reads. The support rows $H,C,P,N$ preserve, replay, source, and name the displayed carrier; they do not add a second window, a second ledger, or an unrecorded Real seal. Hence every public read of $w,r,d,e$ is exhausted by the four sibling surfaces already displayed in the carrier.
\end{proof}

\begin{theorem}[Diagonal-limit representative completion boundary]
\label{thm:diagonal-limit-representative-completion-boundary}
An accepted $\DiagonalLimitRepresentativeUp$ packet supplies the diagonal representative boundary
$$
X,\mu,n,k,i,w,r,d,e,H,C,P,N
$$
and nothing stronger. In particular, it does not supply a host limit, quotient equality of streams, countable choice for a family of representatives, or an ambient $\RealUp$ completeness principle.
\end{theorem}
\begin{proof}
The positive content is exactly the carrier and certificate surface. By \autoref{thm:diagonal-limit-representative-namecert-obligations}, the packet admits the family, modulus, request, computed indices, selected window, readback, ledger, Real seal, transport, continuation, provenance, and local naming rows.

The consumer-facing route is already exhausted by \autoref{thm:diagonal-limit-representative-window-ledger-exhaustion}. Any read of the Real-facing coordinate must pass through $\StreamNameUp$, $\RegSeqRatUp$, $\DyadicRatCoreUp$, and the displayed $\RealUp$ seal row of the same packet.

Finally use the nonchoice boundary of \autoref{thm:diagonal-limit-representative-nonchoice-boundary}. A host limit, quotient stream equality, countable choice principle, or Real-completeness theorem would require coordinates not present in the carrier. The packet therefore supplies only the diagonal representative boundary.
\end{proof}

\begin{theorem}[Diagonal-limit representative tail-selector exclusion]
\label{thm:diagonal-limit-representative-tail-selector-exclusion}
For every accepted $\DiagonalLimitRepresentativeUp$ carrier
$$
D=(X,\mu,n,k,i,w,r,d,e,H,C,P,N),
$$
the displayed tuple excludes host limits, quotient stream equality, hidden tail selectors, and countable choice rows. Any consumer of $D$ can read only the carrier rows, the local $\NameCert_{\DiagonalLimitRepresentativeUp}$ obligations, the real-seal factorization route, the window-ledger exhaustion route, and the completion boundary already displayed by the packet.
\end{theorem}
\begin{proof}
Start from the carrier of \autoref{def:diagonal-limit-representative-carrier}. Its coordinates are exactly the regular-Cauchy family row, modulus row, precision request, selected indices, selected $\StreamNameUp$ window, $\RegSeqRatUp$ readback, $\DyadicRatCoreUp$ ledger, $\RealUp$ seal, and the support rows $H,C,P,N$.

Apply the NameCert obligations of \autoref{thm:diagonal-limit-representative-namecert-obligations}. They account for every admitted row and for the two displayed index computations, so a hidden selector or countable family choice would require an additional certificate coordinate.

Use \autoref{thm:diagonal-limit-representative-real-seal-factorization} and \autoref{thm:diagonal-limit-representative-window-ledger-exhaustion}. These theorems force every Real-facing read through $w,r,d,e$ and through the same support rows, not through a host limit or quotient equality of streams.

Finally apply \autoref{thm:diagonal-limit-representative-completion-boundary}. The boundary theorem states that the packet supplies exactly the displayed diagonal representative boundary and nothing stronger. Therefore the excluded objects are not exports of the accepted carrier.
\end{proof}

\begin{theorem}[Diagonal-limit representative StreamName--RegSeqRat handoff]
\label{thm:diagonal-limit-representative-streamname-regseqrat-handoff}
Every accepted $\DiagonalLimitRepresentativeUp$ consumer route to the seal row $e$ factors first through the selected $\StreamNameUp$ window $w$, then through the selected $\RegSeqRatUp$ readback $r$, and only then through the downstream $\DyadicRatCoreUp$ ledger $d$ before reaching the $\RealUp$ seal row. No route may pass directly from the modulus row $\mu$ or request row $n$ to $\RealUp$ without the $w,r$ handoff.
\end{theorem}
\begin{proof}
Project the accepted carrier. The tuple records $k=\mu(1/(n+1))$ and $i=k+n+1$ before it records the selected window $w$ and selected readback $r$.

Use \autoref{thm:diagonal-limit-representative-namecert-obligations}. The local certificate admits the finite-window exposure through $w$ and the regular-rational readback through $r$ as separate obligations, so the packet cannot collapse them into a direct $\RealUp$ read.

Next apply \autoref{thm:diagonal-limit-representative-real-seal-factorization}. It orders every Real-facing consumer read through $\mu,n,k,i,w,r,d,e,H,C,P,N$, placing $\StreamNameUp$ before $\RegSeqRatUp$ and both before the seal.

The exhaustion theorem \autoref{thm:diagonal-limit-representative-window-ledger-exhaustion} supplies the final constraint. Public reads of the window, readback, ledger, and seal coordinates are exhausted by $w,r,d,e$ with $H,C,P,N$ only transporting, replaying, sourcing, and naming those rows. Thus the route factors through $\StreamNameUp$ and $\RegSeqRatUp$ before it reaches $\RealUp$.
\end{proof}

\begin{theorem}[Diagonal-limit representative dyadic-real seal nonescape]
\label{thm:diagonal-limit-representative-dyadic-real-seal-nonescape}
In an accepted $\DiagonalLimitRepresentativeUp$ carrier, the $\DyadicRatCoreUp$ ledger $d$ and $\RealUp$ seal $e$ are consumed only through the displayed rows
$$
d,e,H,C,P,N.
$$
The carrier exposes no second dyadic ledger, no off-packet real seal, no ambient completeness field, and no route from $d$ or $e$ that bypasses the displayed transport, continuation, provenance, and local naming rows.
\end{theorem}
\begin{proof}
Read the dyadic and real coordinates from \autoref{def:diagonal-limit-representative-carrier}. The only dyadic arithmetic row is $d$, and the only Real-facing seal row is $e$.

Apply the NameCert obligations. They identify dyadic ledger exposure through $d$, terminal seal exposure through $e$, componentwise transport through $H$, continuation replay through $C$, package provenance through $P$, and local naming through $N$ as the complete local surface.

Use \autoref{thm:diagonal-limit-representative-real-seal-factorization} to place $d$ before $e$ in every Real-facing route. Then use \autoref{thm:diagonal-limit-representative-window-ledger-exhaustion} to see that all public reads of $d$ and $e$ are exhausted by the displayed sibling rows and their support.

Finally invoke \autoref{thm:diagonal-limit-representative-completion-boundary}. Since the packet supplies no ambient Real-completeness principle or off-packet seal, consumption of the dyadic ledger and Real seal remains confined to $d,e,H,C,P,N$.
\end{proof}

\begin{theorem}[Diagonal-limit representative terminal uniqueness]
\label{thm:diagonal-limit-representative-terminal-uniqueness}
Let
$$
D=(X,\mu,n,k,i,w,r,d,e,H,C,P,N)
$$
and
$$
D'=(X',\mu',n',k',i',w',r',d',e',H',C',P',N')
$$
be accepted $\DiagonalLimitRepresentativeUp$ carriers. Suppose the displayed modulus, request, selected index, window, readback, dyadic ledger, seal, transport, continuation, provenance, and local naming rows of $D$ and $D'$ are classified together by the componentwise $\hsame$ comparisons and $\Cont$ routes recorded in their packets. Then the two terminal $\RealUp$ seal rows $e$ and $e'$ are classified together by the same representative route, and the accepted classifier exports no second diagonal representative, second terminal seal, hidden tail selector, host limit, quotient stream equality, countable-choice row, or ambient completeness witness.
\end{theorem}
\begin{proof}
Start with the local NameCert obligations. \autoref{thm:diagonal-limit-representative-namecert-obligations} says that every accepted representative packet exposes the modulus request, the two computed indices, the window, readback, dyadic ledger, terminal seal, and the support rows $H,C,P,N$ as the whole local certificate surface.

Next use index stability, \autoref{thm:diagonal-limit-representative-index-stability}. The displayed equations $k=\mu(1/(n+1))$ and $i=k+n+1$ are preserved by the accepted classifier, so the two packets cannot diverge by replacing the selected family index or selected readback index with another selector.

Now apply real-seal factorization and dyadic-real seal nonescape. By \autoref{thm:diagonal-limit-representative-real-seal-factorization}, every Real-facing read reaches the seal only after $\mu,n,k,i,w,r,d$ have been exposed on the same packet. By \autoref{thm:diagonal-limit-representative-dyadic-real-seal-nonescape}, the dyadic and Real coordinates are consumed only through $d,e,H,C,P,N$ and their primed counterparts.

The hypotheses identify exactly those rows on both packets. Therefore the two terminal seal rows are the same accepted representative terminal for the classifier. The nonchoice boundary and completion boundary, \autoref{thm:diagonal-limit-representative-nonchoice-boundary} and \autoref{thm:diagonal-limit-representative-completion-boundary}, exclude a second representative or any off-packet analytic object.
\end{proof}

\begin{theorem}[Diagonal-limit representative carrier handoff]
\label{thm:diagonal-limit-representative-carrier-handoff}
For every accepted $\DiagonalLimitRepresentativeUp$ carrier
$$
D=(X,\mu,n,k,i,w,r,d,e,H,C,P,N),
$$
the finite consumer surface exposes the four sibling rows in the ordered handoff
$$
\StreamNameUp \longrightarrow \RegSeqRatUp \longrightarrow \DyadicRatCoreUp \longrightarrow \RealUp.
$$
The row $w$ is the only $\StreamNameUp$ window, $r$ is the only $\RegSeqRatUp$ readback, $d$ is the only $\DyadicRatCoreUp$ ledger, and $e$ is the only terminal $\RealUp$ seal available to the accepted representative packet.
\end{theorem}
\begin{proof}
Project the carrier of \autoref{def:diagonal-limit-representative-carrier}. Its displayed rows contain exactly one window coordinate $w$, one readback coordinate $r$, one dyadic ledger coordinate $d$, and one Real-facing seal coordinate $e$.

Apply \autoref{thm:diagonal-limit-representative-namecert-obligations}. The local certificate assigns those four coordinates to the sibling surfaces in the stated order and accounts for all support through $H,C,P,N$.

Next use \autoref{thm:diagonal-limit-representative-real-seal-factorization}. Every consumer read that reaches $e$ must pass through $\mu,n,k,i,w,r,d$ first, so the terminal Real row cannot be consumed before the StreamName, RegSeqRat, and DyadicRatCore rows have been read.

Finally apply \autoref{thm:diagonal-limit-representative-window-ledger-exhaustion}. The public reads of $w,r,d,e$ are exhausted by the displayed packet, while the support rows only transport, replay, source, and name those rows. Hence the ordered handoff is the complete finite route exposed by $D$.
\end{proof}

\begin{theorem}[Diagonal-limit representative terminal-seal choice-free totality]
\label{thm:diagonal-limit-representative-terminal-seal-choice-free-totality}
For every accepted $\DiagonalLimitRepresentativeUp$ carrier
$$
D=(X,\mu,n,k,i,w,r,d,e,H,C,P,N),
$$
the terminal seal row $e$ is determined by the finite packet. More precisely, a Real-facing or Completion-facing consumer reaches $e$ only after reading the family row $X$, modulus row $\mu$, request $n$, computed indices $k,i$, window row $w$, readback row $r$, dyadic ledger $d$, and support rows $H,C,P,N$. No host limit, quotient stream, hidden tail selector, selected diagonal family, or countable-choice coordinate can supply an alternate terminal seal.
\end{theorem}
\begin{proof}
Project the accepted packet. The carrier handoff theorem fixes the ordered sibling route through $w,r,d,e$, while the selector exactness theorem fixes $k$ and $i$ as the displayed computations from $\mu$ and $n$.

Next read the seal boundary. The Real-seal route theorem confines every endpoint read to the finite path through $w,r,d,H,C,P,N$, and terminal uniqueness rules out a second seal row once the displayed packet rows are classified together.

Finally apply the nonchoice and completion boundaries. These boundaries exclude hidden selectors, host limits, quotient streams, and countable-choice rows, so the terminal seal is total for the accepted finite packet and cannot be replaced by an ambient analytic object.
\end{proof}

\begin{theorem}[Diagonal-limit representative cofinal-window route exactness]
\label{thm:diagonal-limit-representative-cofinal-window-route-exactness}
For every accepted $\DiagonalLimitRepresentativeUp$ carrier
$$
D=(X,\mu,n,k,i,w,r,d,e,H,C,P,N),
$$
each cofinal refinement that preserves the accepted representative classifier has exactly the route
$$
\mu,n,k,i,w,r,d,e,H,C,P,N.
$$
The refined route is a finite $\StreamNameUp$--$\RegSeqRatUp$--$\DyadicRatCoreUp$--$\RealUp$ route over the same packet, not a detached window ledger, quotient stream comparison, host-limit comparison, hidden tail selector, or countable-choice selection.
\end{theorem}
\begin{proof}
Start with cofinal-window compatibility. It states that a displayed cofinal refinement preserving $\mu,k,i$ and the finite window still follows the ordered sibling route through $w,r,d,e$.

Use request-monotone seal routing to place $\mu,n,k,i$ before the window and readback rows. The modulus-window-readback package then shows that every route to the terminal seal has already consumed $\mu,n,k,i,w,r,d$.

The support rows do not add another endpoint. They transport, replay, source, and name the same finite route, and dyadic-real seal nonescape confines the dyadic and Real rows to $d,e,H,C,P,N$.

Thus the cofinal-window route is exact. The completion boundary and tail-selector exclusion remove detached windows, quotient comparisons, host limits, hidden selectors, and choice rows from the accepted route.
\end{proof}

\begin{theorem}[Diagonal-limit representative cofinal-refinement functoriality]
\label{thm:diagonal-limit-representative-cofinal-refinement-functoriality}
For every accepted $\DiagonalLimitRepresentativeUp$ carrier
$$
D=(X,\mu,n,k,i,w,r,d,e,H,C,P,N),
$$
two displayed cofinal request refinements that preserve the accepted representative classifier compose to the same ordered route
$$
\mu,n,k,i,w,r,d,e,H,C,P,N.
$$
The composed refinement is still a finite $\StreamNameUp$--$\RegSeqRatUp$--$\DyadicRatCoreUp$--$\RealUp$ route through $D$. It exports no second selector, no second window, no quotient-stream comparison, no host limit, no choice row, and no $\RealUp$-completeness coordinate.
\end{theorem}
\begin{proof}
Apply cofinal-window route exactness to the first displayed refinement. Its classifier-preserving route is exactly $\mu,n,k,i,w,r,d,e,H,C,P,N$, with the StreamName, RegSeqRat, DyadicRatCore, and Real rows read inside the accepted packet.

Apply the same exactness theorem to the second displayed refinement after it is viewed over the already refined request. Since both refinements preserve the representative classifier, request-monotone seal routing keeps $\mu,n,k,i$ before $w,r,d,e$ at both stages.

Now compose the two reads. The modulus-window-readback package identifies the same window, readback, dyadic ledger, and Real seal rows, so the second stage has no fresh endpoint to append beyond $w,r,d,e,H,C,P,N$.

Finally use dyadic-real seal nonescape and the completion boundary. They confine the dyadic and Real rows to the displayed seal route and exclude quotient-stream, host-limit, choice, and ambient-completeness coordinates. Hence successive classifier-preserving cofinal refinements collapse to the same finite L10 route through $D$.
\end{proof}

\begin{theorem}[Diagonal-limit representative Real-seal consumer determinacy]
\label{thm:diagonal-limit-representative-real-seal-consumer-determinacy}
For every accepted $\DiagonalLimitRepresentativeUp$ carrier
$$
D=(X,\mu,n,k,i,w,r,d,e,H,C,P,N),
$$
any $\RealUp$ consumer of the representative endpoint is determined by the finite packet route
$$
X,\mu,n,k,i,w,r,d,e,H,C,P,N.
$$
Two accepted consumer reads with the same displayed rows and the same recorded $\hsame$ and $\Cont$ support receive the same terminal seal row $e$ and the same $\StreamNameUp$, $\RegSeqRatUp$, and $\DyadicRatCoreUp$ predecessors. The consumer receives no host limit, quotient stream, hidden tail selector, detached window, or countable-choice row.
\end{theorem}
\begin{proof}
First identify the consumer surface. The dependency-ready package binds the StreamName, RegSeqRat, DyadicRatCore, and Real rows into one local $\NameCert_{\DiagonalLimitRepresentativeUp}$ surface, and the carrier handoff orders those rows as $w,r,d,e$.

Next attach the representative route. Selector exactness fixes $k$ and $i$, request-monotone seal routing fixes the read order before $w,r,d,e$, and terminal-seal choice-free totality fixes the endpoint row $e$ from the same finite packet.

Now compare two accepted reads. If their displayed rows and support routes agree, terminal uniqueness identifies the seal rows, while window-ledger exhaustion identifies the $w,r,d$ predecessors on the same public packet.

The remaining exclusions follow from nonchoice boundary, tail-selector exclusion, and completion boundary. Since none of the excluded objects is a coordinate of $D$, the Real-seal consumer is determined by the finite route alone.
\end{proof}

\begin{theorem}[Diagonal-limit representative selector exactness]
\label{thm:diagonal-limit-representative-selector-exactness}
For every accepted $\DiagonalLimitRepresentativeUp$ carrier
$$
D=(X,\mu,n,k,i,w,r,d,e,H,C,P,N),
$$
the selector coordinates available to the classifier are exactly
$$
k=\mu(1/(n+1))
\quad\text{and}\quad
i=k+n+1.
$$
No hidden tail selector, off-packet family index, countable-choice row, or quotient-stream selector can replace $k$ or $i$ in the accepted classifier.
\end{theorem}
\begin{proof}
Start from index stability, \autoref{thm:diagonal-limit-representative-index-stability}. The displayed equations defining $k$ and $i$ are preserved by accepted classifier transport, so a classified packet cannot change the family index route or the readback index route.

Apply terminal uniqueness, \autoref{thm:diagonal-limit-representative-terminal-uniqueness}. Once the displayed modulus, request, index, window, readback, dyadic, seal, and support rows are classified together, the terminal Real seal is already fixed by the same representative route.

It remains to exclude replacement selectors. The nonchoice boundary and the tail-selector exclusion theorem, \autoref{thm:diagonal-limit-representative-nonchoice-boundary} and \autoref{thm:diagonal-limit-representative-tail-selector-exclusion}, give no packet coordinate for a hidden tail selector, countable choice row, or quotient-stream selector. Therefore the only classifier selector coordinates are $k$ and $i$.
\end{proof}

\begin{theorem}[Diagonal-limit representative Real-seal route]
\label{thm:diagonal-limit-representative-real-seal-route}
Every accepted $\DiagonalLimitRepresentativeUp$ consumer that reaches the terminal row $e$ first reads the sibling rows $w,r,d$ and the support rows $H,C,P,N$. The terminal $\RealUp$ seal is a packet row reached by this finite route, not a $\RealUp$ completeness principle, host limit, quotient of streams, or ambient analytic object.
\end{theorem}
\begin{proof}
Begin with dyadic-real seal nonescape, \autoref{thm:diagonal-limit-representative-dyadic-real-seal-nonescape}. It confines consumption of the dyadic ledger and terminal seal to $d,e,H,C,P,N$.

Use the carrier handoff theorem just proved. Before a consumer can read $d$ or $e$, it must pass through the selected $\StreamNameUp$ window $w$ and the selected $\RegSeqRatUp$ readback $r$ on the same packet.

Next apply \autoref{thm:diagonal-limit-representative-completion-boundary}. The boundary identifies the Real-facing row as the displayed terminal seal of the representative packet and withholds any ambient completeness witness.

Finally invoke terminal uniqueness. Since every accepted route to $e$ has already consumed $w,r,d,H,C,P,N$, and no second terminal seal is exported, the Real-facing endpoint is a finite packet row rather than a Real-completeness principle.
\end{proof}

\begin{theorem}[Diagonal-limit representative cofinal-window compatibility]
\label{thm:diagonal-limit-representative-cofinal-window-compatibility}
Let $D=(X,\mu,n,k,i,w,r,d,e,H,C,P,N)$ be an accepted $\DiagonalLimitRepresentativeUp$ carrier, and let a displayed cofinal refinement of the request preserve the same modulus row $\mu$, the same selected diagonal index route $k,i$, and a refined finite $\StreamNameUp$ window that reads the same $\RegSeqRatUp$ row $r$, $\DyadicRatCoreUp$ ledger $d$, and $\RealUp$ seal row $e$ through the recorded support rows. Then the refined consumer follows the same ordered route
$$
\StreamNameUp \longrightarrow \RegSeqRatUp \longrightarrow \DyadicRatCoreUp \longrightarrow \RealUp
$$
inside $D$. It introduces no hidden selector, no second window ledger, no quotient-stream comparison, no host limit, and no ambient completeness coordinate.
\end{theorem}
\begin{proof}
Project the carrier. The request, selected indices, finite window, regular-rational readback, dyadic ledger, Real seal, and support rows are already part of the accepted tuple of \autoref{def:diagonal-limit-representative-carrier}.

Use the StreamName--RegSeqRat handoff, \autoref{thm:diagonal-limit-representative-streamname-regseqrat-handoff}. Any route to the seal must first expose the selected $\StreamNameUp$ window and then the selected $\RegSeqRatUp$ readback before it may consume the downstream dyadic ledger and Real seal.

Next apply window-ledger exhaustion, \autoref{thm:diagonal-limit-representative-window-ledger-exhaustion}. The public reads of $w,r,d,e$ are exhausted by those four sibling surfaces, with $H,C,P,N$ only transporting, replaying, sourcing, and naming the same rows.

The cofinal refinement preserves the displayed request route and the same $r,d,e$ rows, so it has no separate coordinate from which to read a hidden selector or a second window ledger. The completion boundary then removes quotient-stream, host-limit, and ambient-completeness routes. The refined consumer therefore remains on the ordered $\StreamNameUp$ to $\RegSeqRatUp$ to $\DyadicRatCoreUp$ to $\RealUp$ route of $D$.
\end{proof}

\begin{theorem}[Diagonal-limit representative request-monotone seal route]
\label{thm:diagonal-limit-representative-request-monotone-seal-route}
For every accepted $\DiagonalLimitRepresentativeUp$ carrier
$$
D=(X,\mu,n,k,i,w,r,d,e,H,C,P,N),
$$
any consumer that refines the request while preserving the accepted representative classifier must read the rows $\mu,n,k,i$ before it reads $w,r,d,e$. The only support rows available after this ordered read are $H,C,P,N$, and every $\RealUp$ seal consumption remains confined to the displayed route
$$
\mu,n,k,i,w,r,d,e,H,C,P,N.
$$
\end{theorem}
\begin{proof}
First read the index equations from the packet. By \autoref{thm:diagonal-limit-representative-index-stability}, the equations $k=\mu(1/(n+1))$ and $i=k+n+1$ are part of the accepted classifier surface and cannot be replaced by a request-dependent hidden selector.

Next apply real-seal factorization, \autoref{thm:diagonal-limit-representative-real-seal-factorization}. It orders every Real-facing consumer read through $\mu,n,k,i,w,r,d,e,H,C,P,N$, placing the modulus request and computed indices before the window, readback, ledger, and seal.

Finally apply dyadic-real seal nonescape, \autoref{thm:diagonal-limit-representative-dyadic-real-seal-nonescape}. The dyadic ledger and Real seal are consumed only through $d,e,H,C,P,N$, and there is no second dyadic ledger, off-packet seal, ambient completeness field, or bypass around the displayed support.

Thus a request-refinement consumer that preserves the classifier remains monotone in the displayed request route: it reads $\mu,n,k,i$ first, then $w,r,d,e$, and it has only $H,C,P,N$ as support.
\end{proof}

\begin{theorem}[Diagonal-limit representative dependency-ready package]
\label{thm:diagonal-limit-representative-dependency-ready-package}
For every accepted $\DiagonalLimitRepresentativeUp$ carrier
$$
D=(X,\mu,n,k,i,w,r,d,e,H,C,P,N),
$$
the displayed packet is a single dependency-ready boundary for the diagonal representative route. It consumes the $\StreamNameUp$, $\RegSeqRatUp$, $\DyadicRatCoreUp$, and $\RealUp$ rows in the order
$$
\mu,n,k,i,w,r,d,e,H,C,P,N,
$$
and it packages them as one $\Pkg$-sourced $\NameCert_{\DiagonalLimitRepresentativeUp}$ surface. Every downstream consumer that asks for the diagonal representative endpoint must pass through this packet rather than through a detached stream name, detached regular sequence, detached dyadic ledger, or detached real seal.
\end{theorem}
\begin{proof}
Begin with the carrier obligations of \autoref{thm:diagonal-limit-representative-namecert-obligations}. They place the family row, modulus row, request, computed indices, selected window, readback, dyadic ledger, terminal seal, transport, continuation, package, and local naming row in one finite certificate.

Next apply the StreamName--RegSeqRat handoff of \autoref{thm:diagonal-limit-representative-streamname-regseqrat-handoff}. It forces every route to the terminal seal through $w$ and $r$ before the dyadic ledger and the Real-facing row can be consumed.

Use request-monotone seal routing, \autoref{thm:diagonal-limit-representative-request-monotone-seal-route}, to place the modulus request and computed indices before the sibling rows. Then use dyadic-real seal nonescape, \autoref{thm:diagonal-limit-representative-dyadic-real-seal-nonescape}, to confine the dyadic and Real endpoints to $d,e,H,C,P,N$.

These dependencies exhaust the same carrier. Hence $D$ is not a collection of independent endpoint names; it is one dependency-ready package tying the diagonal representative route to the four consumed rows and to the local $\NameCert_{\DiagonalLimitRepresentativeUp}$ boundary.
\end{proof}

\begin{theorem}[Diagonal-limit representative modulus-window-readback package]
\label{thm:diagonal-limit-representative-modulus-window-readback-package}
In an accepted $\DiagonalLimitRepresentativeUp$ carrier
$$
D=(X,\mu,n,k,i,w,r,d,e,H,C,P,N),
$$
every consumer reaches the terminal $\RealUp$ seal row $e$ only through the finite modulus-window-readback route
$$
\mu,n,k,i,w,r,d.
$$
The rows $H,C,P,N$ may transport, replay, source, and name this route, but they do not supply an alternate route to $e$.
\end{theorem}
\begin{proof}
First read the modulus route from \autoref{thm:diagonal-limit-representative-index-stability}. The selected family index and readback index are the displayed computations attached to $\mu$ and $n$; a consumer cannot replace them by another selector while preserving the accepted classifier.

Next pass to the window and readback rows. By \autoref{thm:diagonal-limit-representative-streamname-regseqrat-handoff}, the route to $e$ must expose the selected $\StreamNameUp$ window $w$ and then the selected $\RegSeqRatUp$ readback $r$ before any Real-facing seal is available.

Now consume the dyadic ledger. The window-ledger exhaustion theorem, \autoref{thm:diagonal-limit-representative-window-ledger-exhaustion}, says that public reads of $w,r,d,e$ are exhausted by those displayed sibling rows, with $H,C,P,N$ only preserving the same surface.

Finally use real-seal factorization, \autoref{thm:diagonal-limit-representative-real-seal-factorization}. It orders every Real-facing read through $\mu,n,k,i,w,r,d,e,H,C,P,N$. Therefore the only path to $e$ has already passed through $\mu,n,k,i,w,r,d$, and the support rows cannot create an alternate endpoint.
\end{proof}

\begin{theorem}[Diagonal-limit representative Real--Completion frontier boundary]
\label{thm:diagonal-limit-representative-real-completion-frontier-boundary}
The package
$$
(X,\mu,n,k,i,w,r,d,e,H,C,P,N)
$$
supports the Real-and-Completion frontier as a named diagonal representative endpoint: it reaches the terminal $\RealUp$ seal row through the dependency-ready route of \autoref{thm:diagonal-limit-representative-dependency-ready-package} and the modulus-window-readback route of \autoref{thm:diagonal-limit-representative-modulus-window-readback-package}. It does not claim Real completeness, quotient stream equality, host limits, selected limits, hidden tail selectors, or countable choice.
\end{theorem}
\begin{proof}
The positive frontier content is the accepted representative package. By \autoref{thm:diagonal-limit-representative-dependency-ready-package}, the packet binds the $\StreamNameUp$, $\RegSeqRatUp$, $\DyadicRatCoreUp$, and $\RealUp$ rows into a single $\Pkg$-sourced $\NameCert_{\DiagonalLimitRepresentativeUp}$ boundary.

The endpoint route is finite. By \autoref{thm:diagonal-limit-representative-modulus-window-readback-package}, every consumer reaches the terminal seal $e$ only through $\mu,n,k,i,w,r,d$, with $H,C,P,N$ serving as support for the same displayed route.

It remains to state the frontier. The completion boundary of \autoref{thm:diagonal-limit-representative-completion-boundary} says that the carrier supplies the diagonal representative boundary and nothing stronger. The tail-selector exclusion of \autoref{thm:diagonal-limit-representative-tail-selector-exclusion} rules out hidden selectors, host limits, quotient stream equality, and countable-choice rows.

Therefore the package is ready for Real-and-Completion consumers that need the named diagonal representative endpoint, while the analytic completion principles and choice-sensitive comparisons remain outside this certificate.
\end{proof}

\begin{theorem}[Diagonal-limit compatibility representative handoff]
\label{thm:diagonal-limit-compatibility-representative-handoff}
Let $K=(R,T,S,D,W,Q,L,H_K,C_K,P_K,N_K)$ be an accepted $\DiagonalLimitCompatibilityUp$ packet whose terminal read has the same four L10 faces as an accepted $\DiagonalLimitRepresentativeUp$ carrier
$$
E=(X,\mu,n,k,i,w,r,d,e,H_E,C_E,P_E,N_E),
$$
namely the same displayed $\StreamNameUp$ window, $\RegSeqRatUp$ readback, $\DyadicRatCoreUp$ ledger, and $\RealUp$ seal rows under the recorded support. Then the terminal read of $K$ factors through the representative carrier $E$ before reaching the terminal Real seal. It inherits the representative nonescape boundary: no quotient stream, host real equality, selected limit, hidden tail selector, choice datum, ambient completeness row, or second terminal seal can enter the handoff.
\end{theorem}
\begin{proof}
First read the compatibility packet through open-phase synchronization. \autoref{thm:diagonal-limit-compatibility-open-phase-exit-synchronization} places the terminal read of $K$ on the four finite exit faces $\StreamNameUp(W)$, $\DyadicRatCoreUp(T)$, $\RegSeqRatUp(Q)$, and $\RealUp(L)$, carried by the displayed support rows.

Next use open-phase four-face readiness, \autoref{thm:diagonal-limit-compatibility-open-phase-four-face-readiness}. It says that readiness of the terminal compatibility read is exactly the displayed four-face surface over its support rows, with no quotient stream, host real equality, selected limit, choice datum, completeness datum, or bridge row.

Now compare those faces with the representative carrier. The hypothesis identifies the compatibility window, readback, ledger, and Real seal with the representative rows $w,r,d,e$. By \autoref{thm:diagonal-limit-representative-window-ledger-exhaustion}, every public read of those four representative coordinates is already exhausted by the same sibling surfaces and the support rows $H_E,C_E,P_E,N_E$.

Finally apply the representative dyadic-real seal nonescape theorem, \autoref{thm:diagonal-limit-representative-dyadic-real-seal-nonescape}. Since the shared dyadic ledger and Real seal are consumed only through the displayed representative support, the compatibility terminal read factors through $E$ before it reaches the seal. The excluded host, quotient, selector, choice, completeness, and second-seal routes are absent from both the compatibility readiness surface and the representative carrier boundary.
\end{proof}
